\scanpage{41}
On peut aussi introduire le groupoïde
\unsure{$C_{1,1}^{\mathrm{ext}\,\prime}$|$C_{1,1}^{\mathrm{cat}\,\prime}$}, déduit du \inserted{groupoïde}
$C_{1,1}^{\mathrm{ext}}$ en \unsure{passant au|prenant un}
quotient par l'\unsure{autom} universel $-\mathrm{id}$
--- correspondant au centre \inserted{$\{\pm1\}$} \unsure{du groupe|des groupes}
\unsure{fondamental|fondamentaux} de ce groupoïde\unsure{.} --- on trouve
un groupoïde connexe équivalent \unsure{can.\textsuperscript{t}}
à celui des torseurs sous \inserted{$\unsure{\mathrm{Gl}|\mathrm{GL}}(2,\mathbb Z)'=$}
$\unsure{\mathrm{Gl}|\mathrm{GL}}(2,\mathbb Z)/(\pm1)$,
\unsure{dt|d'où} on tire maintenant une
\uline{équivalence} de \uline{catégories}
\[
(55)\qquad C_{1,1}^{\mathrm{ext}\,\prime}\xrightarrow{\approx}C_{0,3}
\]
rendant commutatif

(56)\quad
\diagram{[Carré : en haut à gauche $C_{1,\unsure{1|0}}^{\mathrm{ext}\,\prime}$ ; en haut à droite $C_{0,3}$ ; en bas à gauche $\operatorname{Tors}(\unsure{\mathrm{Gl}|\mathrm{GL}}(2,\mathbb Z)')$ ; en bas à droite $\operatorname{Tors}(\unsure{\mathrm{GP}|\mathrm{Gp}}(1,\mathbb Z))$. Deux flèches verticales dirigées vers le bas. Chaque liaison horizontale porte $\approx$ ; leurs pointes sont \unsure{dirigées vers la gauche|non discernables}, plus nettement à gauche pour la liaison inférieure.]},

où $\operatorname{Tors}(\unsure{\mathrm{Gl}|\mathrm{GL}}(2,\mathbb Z)')\xrightarrow{\approx}
\operatorname{Tors}(\unsure{\mathrm{GP}|\mathrm{Gp}}(1,\mathbb Z))$ est
le foncteur tautologique associé à l'iso
\[
(57)\qquad \unsure{\mathrm{Gl}|\mathrm{GL}}(2,\mathbb Z)'\xrightarrow{\approx}\unsure{\mathrm{GP}|\mathrm{Gp}}(1,\mathbb Z).
\]

\scanpage{42}
\mspage{309}
En somme, je suis en train de faire une
paraphrase topologique \unsure{isotopique}, en termes
des groupes fondamentaux à lacets, de
la théorie de \uline{Legendre} des \unsure{relations} entre
courbes elliptiques et courbes rationnelles.
J'aimerais expliciter maintenant l'équivalence
de \uline{Legendre} (55), qui avait
été obtenue grâce à la connaissance des
groupes de Teichmüller idoines relatifs
aux deux membres et d'un isomorphisme
ad-hoc entre ces deux groupes,
directement en termes des groupes
fondamentaux \unsure{de types|des types} $\Pi_{0,3}$ \unsure{ou|et} $\Pi_{1,1}$.

Si $\pi'$ est un groupe à lacets de type
\unsure{$1.1$|$1,1$}, il convient d'abord de noter qu'il
y a dans $\operatorname{Autext}_{\mathrm{lac}}(\pi')$ un élément canonique
d'ordre $2$ --- en fait l'unique élément non trivial
du centre --- noté comme de juste $(-1)_{\pi'}$.
\margin{\uline{NB} \unsure{$(-1)_{\pi'}$|$\operatorname{Aut}_{\pi'}$} conserve l'orientation i.e.
est de \unsure{multiplicateur|multiplicité} $+1$}

\scanpage{43}
Ainsi, on peut considérer $\pi'$ comme muni \inserted{\unsure{can.\textsuperscript{t}}}
d'une \deleted{\unsure{action|autre}} opération extérieure de
$G=\mathbb Z/2\mathbb Z$ sur lui, d'où une extension
canonique
\[
(58)\qquad 1\longrightarrow\pi'\mathrel{\text{---}}\unsure{\mathcal E}\longrightarrow\mathbb Z/2\mathbb Z\unsure{\longrightarrow|\mathrel{\text{---}}}1.
\]

Le modèle topologique nous apprend
qu'il y a exactement trois classes de $\pi'$-
conjugaison de scindages de cette
extension, correspondant aux trois
points d'ordre $2$ d'un groupe topologique
\unsure{$E\simeq\mathbf U\times\mathbf U$} --- les trois points fixes de
$x\mapsto-x$ dans $E\setminus\{0\}$. La théorie des
``\unsure{forages} de trous \unsure{simultanés}'', \unsure{relatif|relative} à une
\unsure{situation} d'un groupe \inserted{fini} d'opérateurs extérieurs
\unsure{d'un|de un} groupe à lacets (théorie qui
\unsure{même} dans le cas discret n'a vraiment
été \unsure{écrit} --- il y a pour le moment des
\unsure{énoncés} conjecturaux) permet alors
\margin{mais ici, on est dans un cas \illegible{2} peut s'en tirer avec
des forages \illegible{1} d'un trou à la fois, \unsure{répétés} successivement\unsure{...}}

\scanpage{44}
\mspage{310}

de construire canoniquement un groupe extérieur \inserted{à lacets $\widetilde{\pi}'$} de type $\Pi_{1,4}$, avec une opération extérieure de $G=\mathbb Z/2\mathbb Z$ sur ce groupe, et une partie $J$ \inserted{$=I\setminus\{i_0\}$} de cardinal $3$ de $I=I(\widetilde{\pi}')$ (qui est de cardinal $4$), $J$ \inserted{stable par $G$ i.e. \unsure{$i_0$ fixe par $G$}}, et un iso. \unsure{du|de} groupe $\widetilde{\pi}'(J)$ \inserted{de type $\Pi_{1,1}$} déduit par ``bouchage des trous en $J$'' avec $\pi'$
\[
(59)\qquad \widetilde{\pi}'(J)\simeq\pi',
\]
iso. qui soit compatible avec les opérations de $G$. Dans le modèle topologique, $\widetilde{\pi}'$ est le groupe fondamental \unsure{cat.|ext.} de \inserted{l'espace} $E\setminus {}_2E$, \inserted{espace} sur lequel $G$ opère librement.

\margin{en fait, $G$ opère \unsure{trivialement} sur $I=I(\widetilde{\pi}')$}

L'extension $\unsure{\widetilde{\mathcal E}}$ de $G$ par $\widetilde{\pi}'$ \unsure{correspondant} à l'action de \deleted{$\unsure{\widetilde{\mathcal E}}$} $G$ sur $\widetilde{\pi}'$
\noindent (60)\quad
\diagram{\begin{tikzcd}[ampersand replacement=\&]
1 \arrow[r, dash] \& \widetilde{\pi}' \arrow[r, dash] \& \unsure{\widetilde{\mathcal E}} \arrow[r, dash] \& \mathbb Z/2\mathbb Z \arrow[r, dash] \& 1
\end{tikzcd}},
\inserted{qui} s'identifie \unsure{donc} au groupe fondamental mixte de $(E\setminus {}_2E,G)$, s'identifie donc \unsure{aussi}-

\scanpage{45}

au groupe fondamental de
\[
(61)\qquad V=(E\setminus {}_2E)/\{\pm1\}
\]
qui est une surface top. de type $(0,4)$.

En d'autres termes, $\unsure{\widetilde{\mathcal E}}$ peut être considéré comme un groupe à lacets de type $(0,4)$, pour une structure à lacets \unsure{facile} à décrire en termes de celle de $\widetilde{\pi}'$, et pour laquelle on a un iso. canonique
\[
(62)\qquad I(\unsure{\widetilde{\mathcal E}})\simeq I(\widetilde{\pi}')/G\simeq I(\widetilde{\pi}').
\]
Ainsi on dispose d'un $i_0\in I(\unsure{\widetilde{\mathcal E}})$, ce qui permet de ``boucher le trou en $i_0$'' pour obtenir un groupe \inserted{effectif} à lacets de type $(0,3)$, soit \deleted{$\unsure{\pi=}$}
\[
(63)\qquad \pi=\mathcal L(\pi')
\]
Reprenant en sens inverse la relation entre $\pi$ et $\pi'$, on peut procéder ainsi: à un groupe \deleted{effectif} \inserted{à lacets} de type $\unsure{(0,3)}$, on associe par

\scanpage{46}
\mspage{311}

le foncteur ``forage de trous'' un groupe à lacets \uline{extérieur} \inserted{$\widetilde{\pi}(\unsure{=\widetilde{\mathcal E}|\in\widetilde{\mathcal E}})$} de type $(0,4)$, muni d'un élément $i_0\in I(\widetilde{\pi})$ tel que \deleted{\illegible{1}}
\[
(64)\qquad I(\widetilde{\pi})\simeq\{i_0\}\amalg\underbrace{I(\pi)}_{\text{cardinal }3}
\]
On a alors ($T$ est le module d'orientation de $\pi$)
\[
(65)\qquad \widetilde{\pi}_{\mathrm{ab}}\simeq T^{I(\widetilde{\pi})}\qquad (\widetilde{\pi}_{\mathrm{ab}})_2\simeq(\mathbb Z/2\mathbb Z)^{I(\widetilde{\pi})}
\]
et on a un hom (ext.)
\[
(66)\qquad \unsure{\widetilde{\pi}}\longrightarrow\mathbb Z/2\mathbb Z
\]
canonique, défini par la \unsure{condition} d'être ``\unsure{ramifié}'' en tous les $i\in I(\widetilde{\pi})$ --- i.e. correspondant à l'hom
\[
(67)\qquad (\widetilde{\pi}_{\mathrm{ab}})_2\xrightarrow{\sim}(\mathbb Z/2\mathbb Z)^{I(\widetilde{\pi})}\longrightarrow\mathbb Z/2\mathbb Z
\]
\uline{somme} des coefficients. Le noyau $\widetilde{\pi}'$ de (66) est un groupe à lacets de type $(1,4)$ --- mais ce n'est pas \unsure{tout} à fait un groupe extérieur --- il n'est pas défini ``\unsure{mod.|aux} \unsure{autom.} int'', \deleted{mais} seulement modulo (en plus) l'opération extérieure de \unsure{symétrie} de $\mathbb Z/2\mathbb Z$


\scanpage{47}

dessus. Si on \unsure{choisit} le groupe $\widetilde{\pi}$ (dans sa ``classe extérieure'') -- on aura en tous cas une suite exacte
\[
(68)\qquad 1\longrightarrow\widetilde{\pi}'\longrightarrow\widetilde{\pi}\mathrel{\text{---}}\mathbb Z/2\mathbb Z\longrightarrow 1
\]
et $\widetilde{\pi}'$ est un groupe à lacets, avec
\[
(69)\qquad I(\widetilde{\pi}')\simeq I(\widetilde{\pi})
\]

En bouchant les trous relatifs à $I(\pi)\subset I(\widetilde{\pi})\unsure{\simeq|=}I(\widetilde{\pi}')$, qui est de cardinal $3$, on trouve un groupe \deleted{extérieur} de type $(1,1)$, avec opération de $\mathbb Z/2\mathbb Z$ dessus \inserted{qui n'est \unsure{autre} que l'op. canonique}. Mais en fonction des indéterminations des choix faits, c'est un peu moins qu'un groupe extérieur de type $(1,1)$ i.e. un objet de $C^{\mathrm{ext}}_{1,1}$ -- c'est un tel \unsure{élément}, ``\uline{modulo} symétrie'' -- i.e. un objet de $C^{\mathrm{ext}\prime}_{1,1}$. Pour avoir un objet de $C^{\mathrm{ext}}_{1,1}$, il faut avoir ``choisi un noyau'' $\widetilde{\pi}'$ pour \unsure{l'hom.} extérieur $\widetilde{\pi}\unsure{\longrightarrow|\mathrel{\text{---}}}\mathbb Z/2\mathbb Z$, en tant que groupe extérieur -- c'est le choix d'un objet

\scanpage{48}
\mspage{312}

dans une $(\mathbb Z/2\mathbb Z)$-gerbe, i.e. \inserted{\unsure{ds}} un groupoïde connexe admettant $\mathbb Z/2\mathbb Z$ comme groupe fondamental.
\margin{\unsure{3}}

\unsure{Modulo} une rédaction plus \unsure{soignée}, il me semble que ces réflexions conceptuelles prouvent bel et bien et sans \unsure{bavure|bavures} qu'on a un isomorphisme canonique (au moins \unsure{mod} autom \unsure{intérieurs|intérieur})
\[
(70)\qquad \mathcal T_{1,1}/\{\pm1\}\xrightarrow{\sim}\unsure{\mathfrak S|\mathcal G}_{0,3}
\]
\margin{en fait, on peut \unsure{le|la} préciser \unsure{mod. autom. int.} provenant \unsure{de|des} $\unsure{\mathfrak S|\mathcal G}^{\unsure{1+|i,+}}_{0,3}\simeq\Pi_{0,3}$}
-- l'indétermination \unsure{mod. autom. intérieurs} provenant \unsure{bien} sûr du choix d'un isomorphisme entre $\mathcal L(\Pi_{1,1})$ et $\Pi_{0,3}$.

Ce fait est tout à fait indépendant, \deleted{de} me semble-t-il, de celui que \unsure{l'hom.} naturel
\[
\mathcal T_{1,1}\xrightarrow{\unsure{\sim}}\unsure{\mathrm{Gl}|\mathrm{GL}}(2,\mathbb Z)
\]
soit un isomorphisme. Je ne vois aucun inconvénient à identifier purement et simplement $\mathcal T_{1,1}$ à $\unsure{\mathrm{Gl}|\mathrm{GL}}(2,\mathbb Z)$ à l'occasion\unsure{,|;} par contre il est plus sujet à caution d'identifier $\unsure{\mathfrak S|\mathcal G}_{0,3}$ à $\mathcal T_{1,1}/\pm1$ \unsure{ou à} $\unsure{\mathrm{Gl}|\mathrm{GL}}(2,\mathbb Z)'$

\scanpage{49}

si ce n'est seulement comme groupe extérieur, à cause de l'indétermination par autom. int. dans (70).

Du point de vue des multiplicités modulaires analytiques, on a un \unsure{morphisme} canonique
\[
(71)\qquad \mathbb M^{\mathrm{an}}_{1,1}\longrightarrow
\underbrace{\unsure{\mathbb U|\mathbb M}^{\mathrm{an}}_{0,\unsure{3}}}_{\inserted{\substack{\text{courbe analytique \unsure{complexe|complète}}\\\text{de type $(0,3)$ ``universelle''}}}}
\simeq \mathbb M^{\prime\mathrm{an}}_{0,4}\simeq(U^{\mathrm{an}}_{0,3},\unsure{\mathfrak S|\mathcal G}_3)
\]
\margin{avec un peu plus de soin, on pourrait inclure dans les \unsure{réflexions} aussi la conjugaison complexe, \unsure{de façon|afin} de n'avoir pas \unsure{ds|dans} (73) à se borner aux sous-groupes $\mathcal T^+$ et $\unsure{\mathfrak S|\mathcal G}^+$ \unsure{en qu.}}
\deleted{dont l'effet sur} qui \unsure{induit} un isomorphisme
\[
(72)\qquad \mathbb M^{\mathrm{an}\prime}_{1,1}\xrightarrow{\sim} M^{\mathrm{an}}_{0,3},
\]
et l'effet de (71) sur les groupes fondamentaux extérieurs n'est autre que l'hom \unsure{ext.|cat.}
\[
(73)\qquad \mathcal T^+_{1,1}\longrightarrow\unsure{\mathfrak S|\mathcal G}^+_{0,3}
\]
induit par (70) -- est-il possible de le préciser de façon utile, en un \unsure{hom.} effectif, non seulement extérieur? Il y a un \unsure{revêt.|revêtement} universel canonique $\widetilde{\mathbb M}^{\mathrm{an}}_{1,1}$ de $\mathbb M^{\mathrm{an}}_{1,1}$, fourni par l'espace \unsure{de|des}

\scanpage{50}
\mspage{313}

Teichmüller (le demi-plan de \unsure{Poincaré} $D$ du \S\ 38) -- relatif à ce revêtement universel on a bel et bien un iso. can.
\[
(74)\qquad \pi_1(\unsure{M|\mathbb M}^{\mathrm{an}}_{1,1})\xrightarrow{\sim}\unsure{\mathrm{Sl}|\mathrm{SL}}(2,\mathbb Z)
\]
(pas seulement extérieur), qui n'est autre d'ailleurs, essentiellement, que (50) -- (restreint au sous-groupe d'indice $2$ $\mathcal T^+_{1,1}$ de $\mathcal T_{1,1}$). On a bel et bien un hom canonique.
\[
(75)\qquad D\longrightarrow U^{\mathrm{an}}_{0,3}\qquad(\simeq \unsure{M|\mathbb M}_{1,1}[2]')
\]
\margin{Pour la définition en \unsure{forme|termes} de \unsure{ce morphisme}, cf \S\quad}
compatible avec un hom. sur les groupes d'opérateurs
\[
(76)\qquad \unsure{\mathrm{Sl}|\mathrm{SL}}(2,\mathbb Z)'\longrightarrow
\underbrace{\unsure{\mathrm{Sl}|\mathrm{SL}}(2,\mathbb Z/2\mathbb Z)\xrightarrow{\sim}\unsure{\mathfrak S|\mathcal G}_3}_{\inserted{\substack{\text{\unsure{cet} iso. \unsure{qui} convient}\\\text{ici est explicité dans}\\\text{\S\quad, p.\quad}}}}
\]
de sorte que $D$ apparaît comme un rev. universel de $U^{\mathrm{an}}_{\unsure{0,3|3}}$, donc aussi de $(U^{\mathrm{an}}_{0,3},\unsure{\mathfrak S|\mathcal G}_3)$, et $\unsure{\mathrm{Sl}|\mathrm{SL}}(2,\mathbb Z)'$ comme le $\pi_1$ mixte relatif à ce rev. universel-ci. Pour \inserted{\unsure{pouvoir}} l'identifier à $\pi_1((U^{\mathrm{an}}_{0,3},\unsure{\mathfrak S|\mathcal G}_3);P=\unsure{-\bar{\jmath}|-j})$, il
